\section{Open problems}\label{sec:open-problems}

The remaining questions concern the interface between recurrence, accessibility and
uniformity, rather than a missing numerical adjustment to the proved thresholds.
The three-state path question is
\[
 F_{\mathcal P,1}(3)\ ?
\]
Equivalently, does one three-state controller solve every finite induced path, or
does the eventual path state complexity become four? The realization theorem decides
whether a prescribed contour-plus-matching cycle can exist. It does not decide whether
one fixed controller has the required state-zero-accessible basin structure on every
path. The unrestricted-grid quantity $F_{\mathcal G,1}(3)$ is also open; branching
invalidates the path phase-contact argument even on a unique recurrent cycle.

The exact minimum universal-host size $u_2$ and its subcubic-tree analogue
$u_{2,\mathrm{tree3}}$ are unknown. The current range
$8\le u_2\le u_{2,\mathrm{tree3}}\le5212$ combines a computer-assisted lower bound
with an explicit construction. It neither proves that $5212$ is optimal nor transfers
the individual-controller threshold seven to the one-host quantifier order.
Let $b_2$ be the minimum cardinality of an arbitrary finite complete test family.
The certified selection gives $b_2\le44$; the packing witness proves optimality only
inside the fixed $144$ candidates, and the global minimum remains open.

Beyond one state there is no exact higher-memory radius law here. The sharp period
formula determines the maximum recurrent period over path geometries, not the
realizability of every weighted walk by a common table. A full higher-state realization
theory, including arbitrary four-state path normal forms, also remains open. These
questions should preserve the distinctions between least tagged period and spatial
period, between actual masks and support masks, and between a constructed cycle and
its accessibility. None is resolved by extending a finite scan horizon.
